% chapters/07/praktikum.tex
\begin{praktikumbox}
Praktikum ini mengikuti Tugas Praktikum 05 pada RPS: pemain mengumpulkan
kepingan dan menghindari rintangan, dengan skor yang bertambah setiap kali
kepingan tersentuh. Praktikum ini sekaligus menutup Blok~I.

\begin{langkah}
  \item \textbf{Mainkan programnya.} Ketik ulang \code{ch07_tabrakan.py} dari
        Daftar~\ref{lst:ch07-tabrakan} ke dalam folder \code{code}. Jalankan,
        lalu mainkan sampai seluruh kepingan terkumpul. Pemain dapat berdiri
        di atas bongkahan dan terhalang oleh sisinya.

  \item \textbf{Periksa penanganan tabrakan secara otomatis.} Jalankan program
        dengan masukan otomatis pada tiga nilai laju bingkai, lalu catat
        jumlah bingkai bertumpang tindih pada masing-masing.
\begin{kodekeluaran}
python ch07_tabrakan.py --frames 300 --auto-run --fps 30
python ch07_tabrakan.py --frames 300 --auto-run --fps 60
python ch07_tabrakan.py --frames 300 --auto-run --fps 120
\end{kodekeluaran}

  \item \textbf{Rusak dengan sengaja, lalu amati.} Satukan kedua sumbu pada
        fungsi \code{update} pemain: pindahkan posisi pada kedua sumbu
        sekaligus, lalu periksa tabrakan sekali saja. Jalankan dan mainkan.
        Catat gejala yang Anda amati, khususnya ketika pemain menabrak sisi
        bongkahan.

  \item \textbf{Rusak dengan cara kedua.} Kembalikan program ke keadaan semula,
        lalu hapus penolan kecepatan tegak sesudah mendarat. Jalankan, lalu
        perhatikan perilaku pemain ketika berdiri di atas bongkahan.

  \item \textbf{Perbaiki kembali, lalu uji kecepatan tinggi.} Kembalikan
        program ke keadaan semula, lalu ubah \code{MOVE_SPEED} menjadi 1500.
        Periksa apakah penghitung tumpang tindih masih bernilai nol. Apabila
        tidak, jelaskan mengapa demikian dengan merujuk
        Bagian~\ref{sec:contoh-bab-7}.

  \item \textbf{Tambahkan satu bongkahan dan dua kepingan} pada tempat yang
        menuntut pemain melompat terlebih dahulu. Pastikan keduanya dapat
        dijangkau.

  \item \textbf{Tulis catatan praktikum} yang memuat ketiga gejala yang Anda
        amati pada langkah ketiga dan keempat, beserta perbaikan yang
        dilakukan. Catatan ini menjadi bukti indikator \ind{c}.

  \item \textbf{Simpan pekerjaan} ke repositori Git dengan pesan penyimpanan
        yang menerangkan isinya.
\end{langkah}
\end{praktikumbox}

\begin{table}[htbp]
\centering
\caption{Lembar kerja Praktikum 05}
\label{tab:lembar-praktikum-05}
{\setlength{\tabcolsep}{4pt}\footnotesize%
\begin{tabularx}{\textwidth}{@{}C{1.8cm} C{3.0cm} C{3.0cm} Y@{}}
\toprule
\headrow \textbf{Batas laju} & \textbf{Bingkai bertumpang tindih} & \textbf{Kepingan terkumpul} & \textbf{Catatan} \\
\midrule
30 FPS & & & \\[1.2em]
\addlinespace
60 FPS & & & \\[1.2em]
\addlinespace
120 FPS & & & \\[1.2em]
\addlinespace
\multicolumn{4}{@{}p{\textwidth}@{}}{\footnotesize\itshape
Gejala ketika kedua sumbu disatukan: \dotfill} \\[1.4em]
\addlinespace
\multicolumn{4}{@{}p{\textwidth}@{}}{\footnotesize\itshape
Gejala ketika penolan kecepatan tegak dihapus: \dotfill} \\[1.4em]
\addlinespace
\multicolumn{4}{@{}p{\textwidth}@{}}{\footnotesize\itshape
Hasil pada kecepatan 1500 piksel per detik: \dotfill} \\[1.4em]
\bottomrule
\end{tabularx}}
\end{table}

\begin{catatabox}
Langkah ketiga dan keempat sengaja meminta Anda merusak program yang sudah
benar. Kebiasaan mengamati gejala dari kesalahan yang dibuat dengan sengaja
akan sangat berguna ketika kesalahan serupa muncul tanpa sengaja pada proyek
akhir nanti: Anda akan mengenali gejalanya lebih cepat karena pernah melihatnya.
\end{catatabox}
